\begin{definition}[The cited results at a setting, carried as a hypothesis package]
Let $\mathbb G$ be a Carnot setting (\noderef{CarnotSetting}) with dimension $n$, number of
generators $m$, step $r$, group law $*$, first-layer projection $\pi_1$, homogeneous norm
$\|\cdot\|$ and horizontal frame $\sigma$; let $d(x,y):=\|y^{-1}*x\|=\|(-y)*x\|$ be the
associated left-invariant distance; let $|E|$ denote Lebesgue measure; and let $\mathcal L$
be the operator (1.8) determined by $(A,H)$. A \emph{cited-results package at
$\mathbb G$} is a Euclidean cited-results package (\noderef{EuclidCitedResults}) together
with a proof of all seven statements below. The setting $\mathbb G$ is fixed throughout: each
of the seven is a statement about $\mathbb G$ itself, and none of them is asserted for any
other setting.

\begin{enumerate}
\item (\textsc{strongMax}; Bardi--Goffi, Corollary Corollary~2.4, sub-solution side.) For every
  $A:\mathbb R^{m}\to\mathbb R^{m\times m}$ and
  $H:\mathbb R^{m}\to\mathbb R$ satisfying the structure conditions
  (\noderef{StructureConditionsES}), every \emph{connected} open
  $U\subseteq\mathbb R^{n}$ and every $u:\mathbb R^{n}\to\mathbb R$ that is a viscosity
  sub-solution of $\mathcal Lu=0$ in $U$ (\noderef{IsViscSubsol}): if $x_0\in U$ satisfies
  $u(y)\le u(x_0)$ for all $y\in U$ and $u(x_0)\ge0$, then $u(y)=u(x_0)$ for all $y\in U$.

\item (\textsc{strongMin}; Corollary Corollary~2.4, super-solution side.) For every
  $A:\mathbb R^{m}\to\mathbb R^{m\times m}$ and
  $H:\mathbb R^{m}\to\mathbb R$ satisfying the structure conditions
  (\noderef{StructureConditionsES}), every \emph{connected} open
  $U\subseteq\mathbb R^{n}$ and every $v:\mathbb R^{n}\to\mathbb R$ that is a viscosity
  super-solution of $\mathcal Lv=0$ in $U$ (\noderef{IsViscSupersol}): if $x_0\in U$
  satisfies $v(x_0)\le v(y)$ for all $y\in U$ and $v(x_0)\le0$, then $v(y)=v(x_0)$ for all
  $y\in U$.

\item (\textsc{pansuDiff}; Pansu's differentiability theorem [P89], as stated in
  Section Section~2.4.) For every open $\Omega\subseteq\mathbb R^{n}$,
  every $f:\mathbb R^{n}\to\mathbb R$ and every $L\in\mathbb R$ such that
  $|f(x)-f(y)|\le L\,d(x,y)$ for all $x,y\in\Omega$: for almost every $x\in\mathbb R^{n}$
  with respect to Lebesgue measure, if $x\in\Omega$ then there is $\xi\in\mathbb R^{m}$ such
  that $f$ is Pansu differentiable at $x$ with horizontal differential $\xi$
  (\noderef{HasPansuGradAt}).

\item (\textsc{pansuConst}; [DMV13, Proposition 4.8].) For every
  \emph{connected} open $\Omega\subseteq\mathbb R^{n}$, every $f:\mathbb R^{n}\to\mathbb R$
  and every $L\in\mathbb R$ such that $|f(x)-f(y)|\le L\,d(x,y)$ for all $x,y\in\Omega$: if
  for almost every $x\in\mathbb R^{n}$ it holds that $x\in\Omega$ implies that $f$ is Pansu
  differentiable at $x$ with horizontal differential $0$, then $f(x)=f(y)$ for all
  $x,y\in\Omega$.

\item (\textsc{wang1}; part (1) of Lemma Lemma~3.8, [CYW07, Proposition 3.3].) For every
  bounded open $\Omega\subseteq\mathbb R^{n}$, every
  $w:\mathbb R^{n}\to\mathbb R$ continuous on $\overline\Omega$ and every $\varepsilon>0$:
  the sup-convolution $w^{\varepsilon}$ (\noderef{SupConv}) is locally semi-convex on
  $\Omega$, and the inf-convolution $w_{\varepsilon}$ (\noderef{InfConv}) is locally
  semi-concave on $\Omega$, i.e. $-w_{\varepsilon}$ is locally semi-convex on $\Omega$.

\item (\textsc{wang2}; part (2) of Lemma Lemma~3.8.) For every bounded
  open $\Omega\subseteq\mathbb R^{n}$ and every $w:\mathbb R^{n}\to\mathbb R$ continuous on
  $\overline\Omega$, the convergences $w^{\varepsilon}\downarrow w$ and
  $w_{\varepsilon}\uparrow w$ hold uniformly on $\overline\Omega$ as
  $\varepsilon\to0^{+}$, in the following three explicit senses:
  \begin{enumerate}
  \item for every $\varepsilon>0$ and every $x\in\overline\Omega$,
    $w(x)\le w^{\varepsilon}(x)$ and $w_{\varepsilon}(x)\le w(x)$;
  \item for all $\varepsilon,\varepsilon'$ with $0<\varepsilon'\le\varepsilon$ and every
    $x\in\overline\Omega$, $w^{\varepsilon'}(x)\le w^{\varepsilon}(x)$ and
    $w_{\varepsilon}(x)\le w_{\varepsilon'}(x)$;
  \item for every $\eta>0$ there is $\varepsilon_0>0$ such that for every
    $\varepsilon\in(0,\varepsilon_0)$ and every $x\in\overline\Omega$,
    $|w^{\varepsilon}(x)-w(x)|\le\eta$ and $|w_{\varepsilon}(x)-w(x)|\le\eta$.
  \end{enumerate}

\item (\textsc{wang3}; part (3) of Lemma Lemma~3.8.) For every $(A,H)$
  satisfying the structure conditions (\noderef{StructureConditionsES}), every bounded open
  $\Omega\subseteq\mathbb R^{n}$, all $u,v:\mathbb R^{n}\to\mathbb R$ continuous on
  $\overline\Omega$, every $R\in\mathbb R$ with $|u|\le R/2$ and $|v|\le R/2$ on
  $\overline\Omega$, and every $\varepsilon>0$: if $u$ is a viscosity sub-solution and $v$ a
  viscosity super-solution of $\mathcal L\cdot=0$ in $\Omega$, then $u^{\varepsilon}$ is a
  viscosity sub-solution and $v_{\varepsilon}$ a viscosity super-solution of
  $\mathcal L\cdot=0$ in the reduced domain $\Omega[\varepsilon]$ (\noderef{ReducedDomain}).
\end{enumerate}

Item 2 of the paper's note for the formalization lists five results as cited from the
literature and \emph{not} to be proved, and prescribes that each be stated once, in the form
the paper uses, and carried as an explicit hypothesis, so that the formalized theorems read
``cited results $\Rightarrow$ Theorem Theorem~1.2 / Theorem Theorem~1.3''. The five are: (a) the
Bardi--Goffi principles, here items 1 and 2; (b) Aleksandrov's theorem and (c) Jensen's
lemma, which are the two items of \noderef{EuclidCitedResults} and are inherited by the
present package; (d) Pansu's theorem together with [DMV13, Prop. 4.8], here items 3 and 4;
and (e) parts (1)--(3) of Lemma~3.8, here items 5--7. No axiom is introduced: the package is a
hypothesis, and every result that uses one of the nine statements takes a package as an
argument, at the setting where the statement is used.

\medskip\noindent\textbf{Why the setting is a parameter.} This is the point at which it would
be easiest to assume too much, and the error is not conservative. Items 1--7 are theorems
about genuine Carnot groups, and their proofs use the dilations and the stratification of the
Lie algebra. Neither is a field of \noderef{CarnotSetting}: that structure records the
exponential coordinates, the polynomial group law, the frame, the properties of the
homogeneous quasi-norm, Hörmander's rank condition (1.6), (2.18), left-invariance of
the horizontal derivatives and the form of the $X_i$ on the first layer --- and nothing
about dilations or a graded decomposition. So for an \emph{arbitrary} instance of
\noderef{CarnotSetting} items 1--7 are not known, and may be false. Asserting them for every
instance at once would therefore risk a package that no proof can ever exhibit, and a
hypothesis package that is unsatisfiable makes every theorem depending on it vacuously true
while every verification lane still passes. Carrying $\mathbb G$ as a parameter removes that
risk at the point where it arises: a consumer assumes the package only at the setting it is
actually working in, and a consumer that works in two settings --- as case (iii) of
Theorem~1.3 does, descending to $\mathbb R^{m}$ --- assumes it separately at each.

Item 3 of the paper's note for the formalization is one-directional in exactly the same way,
and it is worth stating why it does not license more. It permits building the structural facts
of Section Section~2.4 into the definition of the setting ``provided every Carnot group in
the coordinates of Section Section~2.4 satisfies that definition''. The proviso guarantees
that every real Carnot group \emph{is} an instance of \noderef{CarnotSetting}; it says
nothing about every instance being a real Carnot group. It therefore licenses stating the
cited results \emph{for} a given group, which is all items 1--7 do, and does not license
quantifying them over all instances.

\medskip\noindent\textbf{Four points that fix the strength of the package.}

First, items 1 and 2 are Corollary Corollary~2.4, not Theorem Theorem~2.3: the note says the
special case used suffices, and Corollary~2.4 is that case. Its domain hypothesis is retained in
full --- $U$ must be \emph{connected} --- because that is where Corollary~2.4 is stated for a
connected $\Omega$, and the proofs below apply it on balls and on reduced domains. Nothing
comparison-flavoured appears here: Bardi and Goffi also prove a \emph{strong comparison}
principle, which the paper explicitly distinguishes from the comparison principle being
proved (end of Section Section~1.2), and assuming that would make both targets vacuous.
Note also what items 1 and 2 need from the setting even at a fixed $\mathbb G$: Theorem
Theorem~2.3, from which Corollary~2.4 is derived, requires Hörmander's rank condition (1.6)
on $X_1,\ldots,X_m$, and (1.6) is a field of \noderef{CarnotSetting}. Without that field
the package would be unsatisfiable at some instances for a second and independent reason: one
can exhibit a structure meeting all the other conditions in which a single horizontal field
admits no admissible test function at an interior maximum, so that a non-constant function is
vacuously a sub-solution attaining a non-negative interior maximum and items 1--2 fail. The
field must therefore not be removed from the setting.

Second, items 3 and 4 are the two halves of (d), stated for $\mathbb G$'s own quasi-distance
$d$ and horizontal differential; the $d$-Lipschitz hypothesis is the one the consuming nodes
verify, and nothing is asserted about the Pansu differential off the full-measure set.

Third, item 4 retains the connectedness of $\Omega$, as in [DMV13, Proposition 4.8]: a
$d$-Lipschitz function with vanishing horizontal differential is constant on a connected
open set, and on a disconnected one it is only locally constant.

Fourth, item 7 is stated for every $R$ with $|u|,|v|\le R/2$ on $\overline\Omega$ rather
than for the single value $R=2\max\{\|u\|_{L^\infty},\|v\|_{L^\infty}\}$ used in (3.15).
Since $\Omega[\varepsilon]$ shrinks as $R$ grows, quantifying over all admissible $R$ is
weaker than fixing the paper's value, and it is the paper's value that the consuming nodes
will supply. Part (4) of Lemma~3.8 is deliberately absent: the note states it is to be
proved, and it is a separate proof node.

\medskip\noindent\textbf{Satisfiability.} Because the note forbids proving the cited results,
this is the one thing the node cannot establish internally, and the parameter is what keeps
it reachable. For any $\mathbb G$ that is a Carnot group in the coordinates of Section
Section~2.4, items 1--7 are the cited theorems themselves and a package at $\mathbb G$
exists. Two instances where this is used are worth naming. A theorem assuming a package at
the $\mathbb G$ it is about is satisfiable whenever $\mathbb G$ is a real Carnot group, which
is the only case the paper is about. And whenever $\mathbb G$ is a setting whose group law is
addition on $\mathbb R^{n}$ with $n=m$ and whose homogeneous norm is the Euclidean one --- the
Euclidean setting, which is the second place a package is assumed --- a package exists for an
elementary reason: there $d(x,y)=|x-y|$, Pansu differentiability is
ordinary differentiability and item 3 is Rademacher's theorem, items 1 and 2 are the classical
strong maximum and minimum principles, and items 5--7 are the classical sup- and
inf-convolution facts.
\end{definition}
